\chapter{Introduction}\label{chap:01-introduction}

\section{Deux conférences en 2018}\label{sec:intro-conferences}

\paragraph{Vancouver, puis Montréal}
Au printemps 2018, Joelle Pineau, professeure à l'université McGill de Montréal et directrice du
laboratoire montréalais de recherche en intelligence artificielle de Facebook, donne une conférence
invitée à ICLR (\emph{International Conference on Learning Representations}), à Vancouver
\parencite{pineau2018iclr}. La conférence s'intitule \enquote{Reproductibilité, réutilisabilité et
robustesse en apprentissage par renforcement profond}. En décembre, Pineau reprend et augmente son propos à NeurIPS
(\emph{Neural Information Processing Systems}), à Montréal, comme conférence Posner : \enquote{Apprentissage
par renforcement reproductible, réutilisable et robuste} \parencite{pineau2018neurips}. ICLR, NeurIPS et
ICML (\emph{International Conference on Machine Learning}) sont les trois congrès principaux de
l'apprentissage automatique.

\paragraph{La question de départ}
Si une équipe publie un résultat, une autre équipe qui suit la recette obtient-elle le même ? Pineau
part d'un sondage de la revue \emph{Nature} : sur \num{1576} chercheurs de toutes disciplines, plus de 70~\%
disaient avoir déjà échoué à reproduire l'expérience d'un collègue, 52~\% voyaient une crise
\enquote{significative} de la reproductibilité et 38~\% une crise \enquote{légère}
\parencite{baker2016reproducibility}. Lors d'un atelier en décembre 2017, elle a posé la même question à
vingt spécialistes de l'apprentissage automatique : 80~\% y voient une crise, légère ou significative
\parencite{pineau2018iclr}. Elle cite ensuite \textcite{buckheit1995wavelab} : un article scientifique sur
le calcul est seulement la publicité de la recherche. La recherche elle-même est l'environnement
logiciel complet et l'ensemble des instructions qui ont produit les figures.

\paragraph{Un domaine qui grandit, des erreurs qui coûtent}
Pineau rappelle que le RL a cité pendant près de vingt ans un seul succès marquant, le programme de
backgammon TD-Gammon \parencite{tesauro1995tdgammon}, avant AlphaGo \parencite{silver2016alphago}. Vers 2000, un atelier sur le RL
réunissait trente-cinq personnes. En 2016, elle compte environ treize mille articles de RL publiés dans
l'année \parencite{pineau2018iclr}. Elle ajoute que le RL s'approche de problèmes qui comptent : la
conception de traitements médicaux, la gestion des ressources naturelles, des réseaux d'énergie et des
systèmes financiers. Le doute venait aussi des praticiens. En février 2018, Alex Irpan, de Google Brain,
écrivait qu'il répondait non à qui lui demandait si le RL pouvait résoudre son problème. Il pensait avoir raison au moins 70~\% du temps \parencite{irpan2018rlhard}. Pineau projette l'adresse de ce
billet à NeurIPS.

\paragraph{Les seize failles}
Pineau montre ensuite, figure après figure, ce qui manque au RL pour mériter la confiance. Les figures
viennent pour la plupart de l'article de son équipe signé Peter Henderson et ses coauteurs
\parencite{henderson2018matters}. Nous y avons relevé seize failles, numérotées F1 à F16 et rangées en
quatre familles : le matériel nécessaire pour refaire une expérience (code, plan, logiciels), le
réglage des hyperparamètres et des références, les graines aléatoires et la statistique, la portée des
conclusions. La plus célèbre est la faille F11 : dix entraînements du même programme, répartis en deux
groupes de cinq, donnent deux courbes qu'un test déclare différentes, avec $t = -9{,}09$ et
$p = 0{,}0016$. Le chapitre~\ref{chap:02-contexte} décrit les seize failles une à une (tableau~\ref{tab:failles}). Pineau ne s'arrête pas au
constat. Elle propose une liste de contrôle de reproductibilité, que NeurIPS adopte en 2019
\parencite{pineau2021improving}. Elle présente aussi le défi de reproduction d'ICLR 2018, où n'importe
qui pouvait tenter de refaire un article soumis.

\section{La question}\label{sec:intro-question}

Huit ans ont passé. Le RL a changé de taille et de terrain : en 2025 et 2026, plus de la moitié des
articles de RL de notre premier tirage appliquent ce qu'ils appellent RL à de grands modèles de langue.
La question de ce rapport est simple : \emph{a-t-on écouté Joelle Pineau ?} Pour chacune des seize failles,
nous demandons si les articles publiés entre 2018 et 2026 la montrent encore. Nous mesurons aussi ce que coûte la faille
quand on refait les expériences de 2018 avec les algorithmes de 2026. Nous prenons sa liste comme une
grille et nous l'appliquons aux articles publiés, puis à nos propres expériences, avec la même
sévérité.

\section{Deux instruments}\label{sec:intro-instruments}

\paragraph{Le sondage}
Le premier instrument relit des articles avec la grille de lecture de Pineau (critères Q1 à Q8),
complétée par sept critères tirés de sa conférence d'ICLR (X1 à X7) et par une seconde grille de douze
critères pour les failles qu'elle laissait sans mesure (Y1 à Y12). Nous avons lu un premier tirage de 50
articles de RL de 2025 et 2026, parmi les \num{1252} articles de RL d'ICLR 2026, d'ICML 2026 et de NeurIPS
2025, puis 50 articles de RL classique de 2026 (robotique, jeux, contrôle, hors modèles de langue), 50 articles de 2018 et 20 articles par année en 2019, 2020,
2022 et 2024. Des agents IA lisent les articles. Chaque décision repose sur une citation
exacte et sa page. Le verdict compare le RL classique de 2026 à notre propre lecture de 2018, avec les
mêmes grilles et les mêmes lecteurs. Les chapitres~\ref{chap:03-sondage-methode} et~\ref{chap:04-sondage-resultats} décrivent
la méthode et les résultats.

\paragraph{Les expériences}
Le second instrument refait les expériences de 2018 sous un protocole écrit avant le premier
entraînement : quatre programmes (PPO en PyTorch et en JAX, SAC, CrossQ) sur quatre tâches MuJoCo, avec
dix graines aléatoires par programme et par tâche, puis vingt. S'y ajoutent trente essais de PPO et
trente de TRPO sur \task{HalfCheetah}, une évaluation sur des graines aléatoires jamais vues à
l'entraînement et une autre sous une physique modifiée. Un volet illustre les mêmes failles sur des
modèles de langue entraînés avec GRPO. Les chapitres~\ref{chap:05-experiences-methode}, \ref{chap:06-experiences-resultats} et~\ref{chap:07-modeles-de-langue} les décrivent.

\section{Contributions}\label{sec:intro-contributions}

\begin{enumerate}
  \item \textbf{Une mesure à grille égale.} La grille de Pineau, réappliquée aux articles de 2026 et
  calibrée sur notre propre lecture de 50 articles de 2018. À grille égale, quatre pratiques ont
  changé dans le RL classique : le lien vers le code (31~\% puis 74~\%), la sensibilité à un réglage
  montrée (44~\% puis 80~\%), le calcul chiffré (27~\% puis 58~\%), une incertitude rapportée (63~\%
  puis 88~\%). D'autres n'ont pas bougé de façon détectable, dont le test statistique (6~\% puis 12~\%).
  \item \textbf{Une seconde grille préenregistrée.} Douze critères (Y1 à Y12), écrits avec leurs règles
  de décision avant tout codage, qui couvrent la mesure de performance, les tâches, les implémentations,
  les versions, la généralisation, la sensibilité et les outils statistiques récents.
  \item \textbf{Une trajectoire.} Quatre années lues entre 2018 et 2026 : les quatre pratiques montent
  lentement, sans saut juste après la liste de contrôle de 2019.
  \item \textbf{Une batterie de contrôles.} Double lecture à l'aveugle, audit des citations, recherche
  mécanique contre les décisions négatives, provenance publiée de chaque lecteur, coupes par provenance,
  conférence et lecteur, sensibilité au modèle des lecteurs et au filtre de titre, confrontation aux
  listes de contrôle remplies par les auteurs.
  \item \textbf{Les expériences de 2018, refaites.} Deux groupes de cinq essais tirés au hasard ne
  trompent un test de Welch que dans 1,2~\% des cas. Triés, ils le trompent dans 31~\%. Garder les trois meilleurs
  essais sur dix gonfle la moyenne de 33~\%. Le meilleur programme change avec la tâche, sauf à budget
  égal avec vingt graines aléatoires, où CrossQ mène sur les quatre tâches. Une physique modifiée de
  20~\% fait passer le premier au dernier rang sur \task{Hopper}.
  \item \textbf{Une explication du $t$ de 9.} Le script publié par Peter Henderson en 2017 mesure la
  dispersion entre épisodes et ignore celle entre graines aléatoires. Appliqué à nos essais de TRPO, il
  dépasse 9,09 dans 38~\% des tirages au hasard.
  \item \textbf{Un pont vers les modèles de langue.} Un GRPO écrit par nous, sur deux petits modèles
  ouverts, avec la vraie récompense (celle qui vérifie la réponse) et une récompense tirée à pile ou face. Avec cinq graines
  aléatoires, le gain moyen de Qwen avec la vraie récompense vaut $+14$ points, avec un intervalle de
  $-3$ à $+34$ qui contient zéro.
  \item \textbf{Un dépôt qui se vérifie.} Les protocoles datés, les lectures avec leurs citations, les
  résultats des entraînements, les scripts d'analyse et des contrôles qui confrontent chaque nombre du document argumenté (\file{docs/conclusions.md})
  et de l'article de blog aux fichiers de résultats (chapitre~\ref{chap:08-logiciel}).
\end{enumerate}

\section{Guide de lecture}\label{sec:intro-guide}

Le chapitre~\ref{chap:02-contexte} pose le cadre formel du RL, décrit les programmes et les tâches, dresse
l'inventaire des seize failles et situe le travail parmi les travaux voisins. Les chapitres~\ref{chap:03-sondage-methode} et~\ref{chap:04-sondage-resultats} donnent la méthode et les résultats du
sondage. Les chapitres~\ref{chap:05-experiences-methode} et~\ref{chap:06-experiences-resultats} font de même pour les
expériences sur MuJoCo. Le chapitre~\ref{chap:07-modeles-de-langue} décrit le volet GRPO. Le chapitre~\ref{chap:08-logiciel} décrit le logiciel, de l'entraînement aux documents. Le chapitre~\ref{chap:09-ecarts}
liste tous les écarts aux trois protocoles. Le chapitre~\ref{chap:10-limites} rassemble les limites. Le chapitre~\ref{chap:11-conclusion} reprend chaque faille avec son verdict. Les annexes donnent les règles
complètes de la grille (annexe~\ref{chap:A-grille}), des tableaux supplémentaires (annexe~\ref{chap:B-tableaux}),
la liste des articles lus (annexe~\ref{chap:C-articles}) et les environnements logiciels
(annexe~\ref{chap:D-environnements}).

Un lecteur pressé lira le résumé, le tableau~\ref{tab:failles}, la section~\ref{sec:sr-principale} et le chapitre~\ref{chap:11-conclusion}. Un lecteur qui veut refaire un calcul trouvera, dans chaque légende, le
fichier de résultats qui le contient.
